\subsection{Gemeinsame Diskrete Verteilung}

\definition \textbf{Gemeinsame diskrete Verteilung}\\
\smalltext{$X_1,\ldots,X_n$ diskret,$\quad W_i \cleq \N,\quad X_i \in W_i$ fast sicher}
$$
    p := \Bigl(p(x_1,\ldots,x_n)\Bigr)_{x_1\in W_1,\ldots x_n\in W_n} 
$$
$$
    p(x_1,\ldots,x_n) = \P\Bigl[ X_1=x_1,\ldots,X_n=x_n \Bigr]
$$

\theorem $\displaystyle\sum_{x_1\in W_1,\ldots,x_n\in W_n} p(x_1,\ldots,x_n) = 1$

\theorem \textbf{Randverteilung}\\
\smalltext{Die einzelnen Verteilungen $p_X$ lassen sich extrahieren:}\\
\subtext{$\forall z \in W_i$:}
\begin{align*}
 \P[X_i = z] = \underset{x_1,\ldots,z,\ldots,x_n}{\sum} p\Bigl( x_1,\ldots,x_{i-1},z,x_{i+1},\ldots,x_n \Bigr)
\end{align*}
\subtext{$X_1,\ldots,X_n$ diskret,$\quad$ gem. Vert. $p = \Bigl( p(x_1,\dots,x_n) \Bigr)_{x_1 \in W_1,\ldots,x_n \in W_n}$}

{\footnotesize
    \remark Nicht umgekehrt: Aus den Randverteilungen lässt sich nichts über die gem. Vert. schliessen.
}
\subsection{Gemeinsame Stetige Verteilung}

\definition \textbf{Gemeinsame Stetige Verteilung}\\
\smalltext{$X_1,\ldots,X_n: \Omega\to\R,\quad f:\R^n\to\R_+,\quad a_1,\ldots a_n \in \R$}
\begin{align*}
        &\P\Bigl[ X_1\leq a_1,\ldots,X_n\leq a_n \Bigr] \\
        &= \int_{-\infty}^{a_1}\cdots\int_{-\infty}^{a_n} f(x_1,\ldots,x_n)\ dx_n\ldots dx_1    
\end{align*}

\subtext{$f$ heisst \textit{gemeinsame Dichte} von $(X_1,\ldots,X_n)$}

{\footnotesize
    \remark Seien $X_1,\ldots,X_n \overset{\text{i.i.d.}}{\sim} \mathcal{X}$:
    \begin{align*}
        &\P\Bigl[ X_1 \leq x,\ldots, X_n \leq x \Bigr] = \prod_{i=1}^n \P[X_i \leq n] = \Bigl( \P[X \leq n] \Bigr)^n
    \end{align*}
}

\newpage

\theorem $\displaystyle\int_{-\infty}^\infty\cdots\int_{-\infty}^\infty f(x_1,\ldots,x_n)\ dx_n\ dx_1 = 1$\\
\subtext{Umgekehrt existiert für jedes solches $f$ ein Raum $(\Omega, \F, \P)$}

{\scriptsize
    \textbf{Beispiel:} Finde $c$ s.d. $f_{X,Y}$ eine Dichte ist, wobei
    $
        f_{X,Y}=\begin{cases}
            ce^{-x}     & 0 < x < y \\
            0           & \text{sonst}
        \end{cases}
    $
    \begin{align*}
        \int_0^\infty \int_0^x ce^{-x}\ dy\ dx  &= c\cdot \int_0^\infty xe^{-x}\ dx                                                     \\ 
                                                &= c\cdot\biggl( \Bigl[ -xe^{-x} \Bigr]_0^\infty + \int_0^\infty e^{-x}\ dx \biggr)     \\
                                                &= c\cdot\biggl( 0 + \Bigl[ -e^{-x} \Bigr]_0^\infty \biggr)                             \\
                                                &= c
    \end{align*}
    Also gilt wegen dem vorherigen Satz: $c=1$.
}

{\footnotesize
    \remark Randdichten: Integration über übrige Variablen.
}

% Bedingte Dichten fehlen noch, siehe HW5, F3

\theorem \textbf{Erwartungswert}\\
\subtext{z.B. $\E[XY] = \int_{-\infty}^{\infty}(xy)\cdot f_{X,Y}(x,y)$.}
\begin{align*}
    &\E\Bigl[ \phi(X_1,\ldots,X_n) \Bigr] \\
    &= \int_{-\infty}^\infty\cdots\int_{-\infty}^\infty\phi(x_1,\ldots,x_n)\cdot f(x_1\cdots x_n)\ dx_n\ldots dx_1    
\end{align*}

% Randverteilungen

\theorem \textbf{Unabhängigkeit}
$$
    X_1,\ldots,X_n \text{ stetig, } f(x_1,\ldots,x_n) = f_1(x_1)\cdots f_n(x_n)
$$

{\scriptsize
    \remark Unabhängige stetige Variablen sind automatisch gemeinsam stetig

    \remark Für unabh. stetige Variablen muss der Träger ein karthesischs Produkt sein. z.B. für 2 Variabeln in $\R$ ginge $[0,1]\times[0,2]$ (Ein Rechteck) 
}

% It's important to him to know \E[X]\E[Y] = \E[X\cdot Y] isn't enough for indep., requires \E[\phi(X)]\E[\psi(X)] instead \forall \phi,\psi

\notation \textbf{Bedingte gem. Verteilung}\\
\subtext{Analog zur Definition von $\P[X\sep Y]$.}
$$
    f_{X \vert Y} := \frac{f_{X,Y}(x,y)}{f_Y(y)}
$$

{\footnotesize
    \textbf{Beispiel}: Bei $f_{X,Y}=e^{-x},\ f_Y(y)=e^{-y}$ erhält man:
    $$
        f_{X|Y} \overset{\text{def.}}{=} \frac{f_{X,Y}(x,y)}{f_Y(y)} = \frac{e^{-x}}{e^{-y}} = e^{-(x-y)}
    $$
    Setzt man $z:=x-y$ gilt $f_{X|Y}\sim\text{Exp}(y)$. Gegeben $\{Y=y\}$ gilt:
    $$
        X - y \sim \text{Exp}(1) \qquad \text{oder} \qquad X \sep Y=y \sim y + \text{Exp}(1)
    $$
}